\chapter{Conclusion}\label{chap:9}

\begin{quote}
\emph{``In theory, theory and practice are the same. In practice, they
are not.''}\\
--- attributed to various sources
\end{quote}

Four ideas run through this manuscript. \emph{Notation is a tool, not
gatekeeping} --- mathematical notation gives a problem a precise shape
that prose cannot, and reading it is a habit anyone working with
optimization can pick up rather than a barrier put up by a community.
\emph{Optimization models decompose into five parts}: the entities the
problem talks about, the parameters that describe an instance, the
decision variables, the constraints that restrict their combinations,
and the objective that ranks the feasible ones. \emph{The verifier is
orthogonal} to whichever solver eventually finds a solution --- it does
not help find one, but it pins the problem definition down in code and
flags any candidate that breaks a rule, which protects every later
modeling decision from drifting away from what was actually meant.
\emph{Knowing your solver's language matters} --- CP-SAT has an
expressive vocabulary (reified Booleans,
\lstinline!add_circuit!,
\lstinline!add_no_overlap!,
\lstinline!add_cumulative!, and the rest of \Cref{chap:7}),
and an idiomatic CP-SAT model often looks different from the same
problem expressed in a MIP solver or another constraint framework, even
when the paper formulation is identical.

Two further moves come up often enough in practice to be worth naming,
even if a thorough treatment is outside this manuscript's scope.
\emph{Trimming the model's scope:} not every dimension a problem carries
needs to be a decision variable. Some ``decisions'' are downstream ---
they follow trivially from the genuinely hard ones and can be resolved
by a post-processing step. An exam planner that fixes day and room and
then packs the times greedily within each day is a typical example.
Whenever a dimension is separable from the rest in this way, pushing it
out of the solver shrinks the search space without changing the answer;
the discipline is verifying that the decomposition really is lossless
before committing to it. \emph{Two models, one verifier:} a model that
mirrors the paper formulation closely is sometimes too slow for the
instances a project actually faces. Once the speed modifications grow
beyond a tweak --- alternative representations, symmetry breaking,
problem-specific reductions --- it can be cleaner to maintain a
\emph{faithful} model alongside a \emph{tractable} one and let the
verifier judge both on the same instances. Most projects never need
this, and benchmarking should drive the decision rather than instinct.

CP-SAT itself covers a large slice of the discrete, combinatorial
landscape, but not all of it: continuous variables push toward LP/MIP
solvers (Gurobi, HiGHS, CPLEX), instances too large for exact methods
push toward metaheuristics, and gradient-based optimization lives in
another world entirely. What carries across is the part of this
manuscript that is not CP-SAT-specific --- the paper notation of \Cref{chap:2} and the coded verifier of \Cref{chap:4}. A
\lstinline!Solution! and an
\lstinline!evaluate(instance, sol)! know nothing about
which solver produced the input, and the habit of writing a yardstick
cheaper than the artifact it judges transfers wherever optimization is
done.
